\section{Calibration-marginal adaptive testing}
\label{sec:framework}

\paragraph{Items, oracles and verdicts.} A test item $i$ is a task with an oracle $O_i$ that
maps a system's output to a verdict $y_i\in\{0,1\}$: exact numeric match, hidden unit tests,
or a log-likelihood comparison of answer options. The framework (EOST, for entropy-oracled
information-optimal self-testing, optimal in the one-step sense of Section~\ref{sec:theory})
has one structural rule: the system under evaluation may propose items but is never the
oracle of its own verdicts, which avoids the circularity of models judging models
\citep{zheng2023judging}. All items here come from fixed banks; self-generated items were
explored only in a historical pilot and are not part of the evidence. Verdicts are
deterministic (one recorded evaluation per system and item) or recorded aggregates (a pass
rate over samples binarized at one half). Figure~\ref{fig:framework} summarizes the test.

\begin{figure}[t]
\centering
\begin{tikzpicture}[
  font=\footnotesize,
  box/.style={draw=black!55, rounded corners=2pt, align=center, inner sep=4pt, fill=white, text width=3.15cm},
  key/.style={draw=blue!60!black, rounded corners=2pt, align=center, inner sep=4pt, fill=blue!6, text width=3.75cm},
  arr/.style={-{Stealth[length=2mm]}, black!70, thick},
  lab/.style={font=\scriptsize, text=black!65, fill=white, inner sep=1pt}
]
\node[box] (bank) at (0,0) {Item source\\{\scriptsize bank, generator, or the evaluated model}};
\node[key] (sel) at (5.5,0) {Select item $i$ maximizing\\ $I(\Theta;y_i\mid\mathcal D_t)$ {\scriptsize (or $R_i$)}};
\node[box] (sys) at (11.0,0) {System under test\\{\scriptsize answers item $i$}};
\node[box] (panel) at (0,-2.1) {Calibration panel $\mathcal C$\\{\scriptsize $N$ systems, target's family or cluster excluded}};
\node[key] (cal) at (5.5,-2.1) {Empirical-Bayes calibration\\{\scriptsize item posteriors $\pi_i(\log a,b)$, ERF}\\{\scriptsize $\bar P_i(\theta)=\mathbb E_{\pi_i}\,\sigma(a(\theta-b))$}};
\node[box, fill=black!4] (oracle) at (11.0,-2.1) {Independent oracle $O_i$\\{\scriptsize unit tests, exact match}};
\node[box] (post) at (11.0,-4.0) {Ability posterior $\pi_t(\theta)$\\{\scriptsize score $\hat A_t$ and 90\% interval}};
\draw[arr] (bank) -- node[lab, above=1pt]{candidates} (sel);
\draw[arr] (panel) -- node[lab, above=1pt]{verdicts} (cal);
\draw[arr] (cal) -- node[lab, right=2pt]{$\bar P_i$} (sel);
\draw[arr] (sel) -- node[lab, above=1pt]{item $i$} (sys);
\draw[arr] (sys) -- node[lab, right=2pt]{output} (oracle);
\draw[arr] (oracle) -- node[lab, right=2pt]{$y_i\in\{0,1\}$} (post);
\draw[arr] (post.west) -- (8.35,-4.0) -- node[lab, pos=0.15, left=1pt, fill=none]{update $\mathcal D_{t+1}$} (8.35,-1.05) -- (6.95,-1.05) -- (6.95,0 |- sel.south);
\end{tikzpicture}
\caption{EOST-M: calibration-marginal adaptive testing with an independent oracle.}
\label{fig:framework}
\end{figure}

\paragraph{Measurement model and calibration.} Let target $s$ have competence $\theta_s$ and
respond according to the two-parameter logistic (2PL) model
\begin{equation}
\Pr(y_{si}=1\mid\theta_s,a_i,b_i)=\sigma\big(a_i(\theta_s-b_i)\big),\qquad a_i>0 .
\end{equation}
Item parameters $\psi_i=(\log a_i,b_i)$ are calibrated on a \emph{panel} $\mathcal C$ of $N$
other systems that excludes the target and every related system (its heuristic family in
study~1, its strict cluster in study~2; Section~\ref{sec:design}). Calibration has four
steps shared by every policy: (i) joint maximum a posteriori (MAP) estimation of the
calibrators' abilities $\hat\theta_{\mathcal C}$ and the item parameters, with
$\theta_j\sim\mathcal N(0,1)$, $\log a_i\sim\mathcal N(0,1)$, $b_i\sim\mathcal N(0,2^2)$;
(ii) an empirical-Bayes hyperprior $\log a_i\sim\mathcal N(\mu_a,\sigma_a^2)$,
$b_i\sim\mathcal N(\mu_b,\sigma_b^2)$ that maximizes the marginal likelihood of all items
given $\hat\theta_{\mathcal C}$ (EM on a $14\times65$ grid), well determined even for $N=4$
because it pools hundreds of items; (iii) plug-in parameters $(\hat a_i,\hat b_i)$, the
conditional MAP under that hyperprior; (iv) the exact grid posterior
$\pi_i(\psi)\propto p(\psi)\prod_{j\in\mathcal C}p(y_{ji}\mid\hat\theta_j,\psi)$ on
$26\times129$ points of $(\log a,b)\in[-2.5,2.5]\times[-8,8]$. The \textbf{plug-in} model
uses $P_i(\theta)=\sigma(\hat a_i(\theta-\hat b_i))$, as in classical CAT and recent LLM work
\citep{hofmann2025fluid,li2026atlas}; the \textbf{marginal} model uses the expected
response function (ERF)
\begin{equation}\label{eq:erf}
\bar P_i(\theta)=\int\sigma\big(a(\theta-b)\big)\,\pi_i(d\psi)
\end{equation}
\citep{mislevy1994erf}, which is the exact response probability given
$\hat\theta_{\mathcal C}$ and the hyperprior (Lemma~\ref{lem:exact}). The two models share
$\hat\theta_{\mathcal C}$ and the hyperprior and differ only in plugging in a mode versus
integrating the item posterior.

\paragraph{Item selection.} The target's ability posterior $\pi_t$ is kept on a 201-point
grid over $[-5,5]$ with prior $\mathcal N(0,1)$. At step $t$ an item is chosen by
\textsc{Random} (uniform choice), \textsc{Entropy} (largest predictive entropy),
\textsc{MFI} (largest Fisher information at the posterior mean), \textsc{PWFI}
(posterior-weighted Fisher information; \citealp{vanderlinden1998bayesian}), \textsc{EIG}
(largest expected information gain $I(\theta;y_i\mid\mathcal D_t)$) or \textsc{EVR}
(largest expected reduction of the posterior variance of the unobserved bank total,
Proposition~\ref{prop:evr}), each under the plug-in (\texttt{:pi}) or the marginal
(\texttt{:mg}) model (PWFI under the marginal model only). The proposed method,
\textbf{EOST-M}, is \textsc{EIG:mg}; the state of practice is \textsc{MFI:pi}; the
model-free baseline \textsc{Naive} samples items uniformly and reports the sample mean with
a Wilson interval and finite-population correction. Entropy, EIG and PWFI are matrix--vector
products with precomputed item-by-grid matrices, EVR uses the closed form of
Proposition~\ref{prop:evr} on the grid, and ties are broken at random.

\paragraph{Score, intervals and recalibration.} The full-bank accuracy
$A=\frac1I\sum_iy_{si}$ is estimated by
$\hat A_t=\frac1I\big(\sum_{\text{asked}}y_{si}+\sum_{\text{not asked}}\mathbb E_{\pi_t}P_i(\theta)\big)$.
Its 90\% interval uses the exact mean and variance of the predictive distribution of the
unobserved total---a mixture over the ability grid of Poisson-binomial sums---with a normal
approximation. Because a model-based interval ignores misspecification and dependence
between systems, we also study a leave-one-calibrator-out recalibration (``+R''): for every
calibrator $c$, the item posteriors are recomputed without $c$'s cluster, the same test is
replayed on $c$, whose score is known, and the recalibrated half-width at step $t$ is the
90\% quantile of the absolute errors pooled over calibrators and over the steps within a
factor 1.5 of $t$. With $N\le8$ pseudo-targets this is a cross-validated correction, not a
conformal guarantee \citep{vovk2005algorithmic}; the rule was chosen among three variants on
development banks before study~2.

\paragraph{Endpoints.} Competence is a model construct whose scale is fixed only by the
prior, so the endpoints are observable. (1) \emph{Held-out predictive log loss}: a holdout of
20\% of the bank, drawn once per target and replicate, never administered and shared by all
policies and panel sizes, is predicted after each step with $\mathbb E_{\pi_t}P_i(\theta)$,
and the log loss is averaged over the steps $t\le T$; this strictly proper score rewards both
accuracy and calibration. (2) \emph{Score error}: the integrated absolute error (IAE), the
average of $|\hat A_t-A|$ over $t\le T$. (3) \emph{Coverage} of the 90\% interval and the
\emph{interval score} $(u-l)+\frac2\alpha(l-A)_++\frac2\alpha(A-u)_+$ with $\alpha=0.1$,
which penalizes both width and misses \citep{gneiting2007strictly}. Secondary endpoints
include the Brier score, the number of items after which the score error stays below a
tolerance, and, on synthetic banks, the error and coverage of the true $\theta$.
